% Generated draft: requires mathematical and source review.
\documentclass[11pt]{article}
\usepackage[margin=25mm]{geometry}
\usepackage{fontspec}
\setmainfont{Latin Modern Roman}
\usepackage{amsmath,amssymb,mathtools,unicode-math}
\setmathfont{Latin Modern Math}
\usepackage{enumitem,needspace,xurl,hyperref}
\setlength{\parindent}{0pt}
\setlength{\parskip}{5pt}
\setlength{\emergencystretch}{3em}
\newcommand{\N}{\mathbb{N}}
\newcommand{\Z}{\mathbb{Z}}
\newcommand{\Q}{\mathbb{Q}}
\newcommand{\R}{\mathbb{R}}
\newcommand{\C}{\mathbb{C}}
\newcommand{\F}{\mathbb{F}}
\newcommand{\A}{\mathbb{A}}
\newcommand{\PP}{\mathbb{P}}
\newcommand{\E}{\mathbb{E}}
\newcommand{\eps}{\varepsilon}
\newcommand{\dd}{\,\mathrm{d}}
\begin{document}
\section*{UnsolvedMath ID: 30002689}
\textbf{Source problem code:} catalogue entry
\subsection{Definitions and mathematical statement}
A Riemannian metric \(g\) on a smooth manifold \(M\) is Einstein if there exists a constant \(\lambda\in\R\) such that
\[
\operatorname{Ric}_g=\lambda g.
\]
It is homogeneous if a Lie group acts transitively on \(M\) by isometries. For a specified transitive action of a Lie group \(H\) on \(M=H/L\), with \(L\) closed, write
\[
\mathcal E_H(H/L)=\{g:\ g\text{ is an }H\text{-invariant Riemannian metric on }H/L,\ \operatorname{Ric}_g=\lambda g\text{ for some constant }\lambda\in\R\}.
\]
Such invariant metrics correspond to positive-definite inner products on \(T_{eL}(H/L)\cong\mathfrak h/\mathfrak l\) preserved by the isotropy representation of \(L\).

Under the usual compact convention, a generalized flag manifold is a quotient \(G/K\), where \(G\) is a compact connected semisimple Lie group and \(K=C_G(T)\) is the centralizer of a torus \(T\subseteq G\). The record does not explicitly state its convention for generalized flag manifolds. The Stiefel spaces are presented in the record as
\[
V_k\R^n=\operatorname{SO}(n)/\operatorname{SO}(n-k),
\]
where the denominator has its standard block embedding. The intended integer range of \(n,k\) is not supplied. In the usual range \(1\leq k<n\), this quotient parametrizes ordered orthonormal \(k\)-frames in \(\R^n\). At \(k=n\), the displayed quotient gives \(\operatorname{SO}(n)\), whereas the space of all ordered orthonormal \(n\)-frames is \(\operatorname{O}(n)\); the intended endpoint convention therefore requires clarification.

An isotropy representation has pairwise nonequivalent irreducible components if, as a real representation of the isotropy group \(L\), it admits a decomposition
\[
\mathfrak h/\mathfrak l=\mathfrak m_1\oplus\cdots\oplus\mathfrak m_r,
\qquad
\mathfrak m_i\text{ irreducible},\qquad
\mathfrak m_i\not\cong\mathfrak m_j\quad(i\ne j).
\]
This uses real irreducibility; the source does not explicitly specify the representation-field convention.

\textbf{Requested classification.} For every generalized flag manifold \(G/K\) and every Stiefel manifold in the intended parameter range, classify all homogeneous Einstein metrics. For the displayed transitive actions, the corresponding invariant-metric formulation is to determine \(\mathcal E_G(G/K)\) and \(\mathcal E_{\operatorname{SO}(n)}(V_k\R^n)\). The record does not resolve whether these fixed-action invariant metrics exhaust the intended classification target, or whether metrics homogeneous under other transitive actions must also be included.

\textbf{Requested finiteness questions.} Determine, for each homogeneous space whose isotropy representation has pairwise nonequivalent irreducible components, whether its homogeneous Einstein metrics form a finite set of classes under the intended equivalence relation. Determine the analogous finiteness question for each compact semisimple Lie group. If the group itself is regarded as a homogeneous space under left translations, its invariant metrics are precisely its left-invariant metrics; the record does not explicitly restrict the group question to this action.

The equivalence relation in both classification and counting is unspecified. In particular, a finiteness question must address constant rescaling, since \(g\mapsto cg\), \(c>0\), preserves the Einstein property. One possible convention is to identify \(g_1,g_2\) when \(g_2=c\phi^*g_1\) for a diffeomorphism \(\phi\) and \(c>0\), but the record does not choose this convention, an equivariant version, or a normalization. Compactness is not explicitly imposed on the general homogeneous spaces in the finiteness question. These unresolved choices prevent a unique fully specified counting statement.
\subsection{Short English statement}
Classify homogeneous Einstein metrics on generalized flag manifolds and the displayed real Stiefel quotients, and determine finiteness for spaces with pairwise inequivalent irreducible isotropy components and for compact semisimple Lie groups. The acting groups, equivalence relation, scaling convention, general compactness assumptions, and Stiefel parameter boundaries require clarification.
\subsection{Sources}
\begin{itemize}
\item Supplied local registry record, permanent UnsolvedMath ID \(30002689\), title \emph{Homogeneous Einstein Metrics on Flag and Stiefel Manifolds}, source problem code \texttt{catalogue entry}, in \emph{Oberwolfach Reports Open Problems}. The record gives proposed year 2014 and names no proposer.
\item Supplied source citation: \emph{Mini-Workshop: Einstein Metrics, Ricci Solitons and Ricci Flow under Symmetry Assumptions} (2015), \emph{Oberwolfach Reports}, pp. 2529--2568, DOI \texttt{10.4171/owr/2014/45}. Supplied source URL and literature reference \emph{Oberwolfach Report 45/2014}: \url{https://doi.org/10.4171/owr/2014/45}. The citation's 2015 date and the report's 2014 identification are reproduced as supplied; their relationship has not been independently checked.
\item Supplied literature reference: \emph{Classification of Einstein Metrics on Flag Manifolds}, \url{https://doi.org/10.1016/j.bulsci.2013.11.002}.
\item Supplied literature reference: \emph{Flag Manifolds with Five Isotropy Summands}, \url{https://doi.org/10.48550/arXiv.1207.2897}.
\item Supplied literature assessment, checked at \texttt{2026-08-21}: the registry labels the problem \texttt{partially\_solved} and states that many generalized flag-manifold families and low-isotropy cases have been classified, while the requested all-metrics classifications for arbitrary flag manifolds, Stiefel manifolds, and compact simple groups remain open. This is supplied dated curation derived from the OWR record, not an independently verified literature assessment. Its reference to compact simple groups does not replace the statement's compact semisimple group target.
\end{itemize}
\end{document}
